\documentclass[12pt]{article}
\usepackage[utf8]{inputenc}
\usepackage[T1]{fontenc}
\usepackage{hyperref}
\title{Chimalapa Zoque Language Resources for K-12 Curriculum Development}
\author{Compiled by Claude Sonnet 5.5 (claude-sonnet-5-5) for Andr\'as Kornai}
\date{}
\begin{document}
\maketitle

\section{Introduction}
Chimalapa Zoque (ISO 639-3 zoh) is an endangered Mixe-Zoquean language of the Chimalapas region in Oaxaca, Mexico. Speaker counts found in the searches vary, and the prompt's figure is about 4,500. The community calls its language \emph{Anepaan}, according to an encyclopedic entry seen in search. The searches found a documentation collection and an older linguistic volume but no school materials.

\section{Primary Educational Resources}
No primers, textbooks or school programmes specific to Chimalapa Zoque were found. The search returned Mexican SEP ``Nueva Escuela Mexicana'' class sheets, but these were not confirmed to concern this language and are not listed as resources.

\section{Linguistic and Reference Materials}
\begin{itemize}
\item \emph{Zoque de Chimalapa}, Archivo de Lenguas Ind\'igenas de M\'exico (collected by L. Knudson), Colegio de M\'exico repository: \url{https://repositorio.colmex.mx/pdf_viewer/fn1072420?locale=en}
\item Endangered Languages Project, resources page for Chimalapa Zoque: \url{https://www.endangeredlanguages.com/es/node/37547/resources}
\end{itemize}

\section{Religious and Community Resources}
\begin{itemize}
\item Redalyc article on territory and identity of the Chimalapa Zoque: \url{https://www.redalyc.org/pdf/4559/455944927004.pdf}
\end{itemize}

\section{Digital and Audio Resources}
\begin{itemize}
\item ELDP archive, ``Documentation of five Zoquean languages spoken in Mexico: Santa Mar\'ia Chimalapa'' (audio and video narratives): \url{https://eldp.access.preservica.com/dk0528/}
\end{itemize}

\section{Curriculum Development Context}
The language is endangered and heavily displaced by Spanish, and no school use is known. The ELDP recordings and the older Colegio de M\'exico volume could support material development. A curriculum would need to be built with Chimalapa speakers, who should also decide how the recorded materials are used.

\end{document}